\subsection{Subcalibrations from orthogonal quartic data}
\label{sec:quartic-table-existence}

The orthogonal-quartic table determines the minimal level, singular
set, and subcalibration of every realization in the following sufficient
minimizing range. Exact thresholds for particular families also use
the instability arguments below.

Throughout this subsection, let $E_0,E_1,E_2$ be nonzero homogeneous
quartic polynomials on $\R^N$ satisfying
\eqref{eq:orthogonal-quartic-data}, with $R=|x|^2$.  Suppose
$r_1,r_2\geq2$, and put
\begin{gather}
 \begin{gathered}
 h=r_1+r_2,\qquad \rho=r_0+r_1+r_2,\\
 c=\frac{r_1-1}{h-2},\qquad a=1-2c,\qquad b=c(1-c)>0,
 \end{gathered}
 \label{eq:quartic-general-weights}\\
 T=\frac{E_1+E_2}{4},\qquad
 F=\frac{(1-c)E_1-cE_2}{4}.
 \label{eq:quartic-general-generators}
\end{gather}
By \cref{lem:quartic-block-weights-image}, $N=2\rho-2$.
The choice of $c$ is the minimal-level condition
\eqref{eq:quartic-level-condition}.  Direct substitution gives
\begin{equation}\label{bh:reduced-table}
 \begin{gathered}
 (\Gamma(I_i,I_j))_{I=(R,T,F)}=
 \begin{pmatrix}
 4R&8T&8F\\
 8T&4RT&4RF\\
 8F&4RF&4R(aF+bT)
 \end{pmatrix},\\
 (\Delta R,\Delta T,\Delta F)=(2N,2hR,2aR).
 \end{gathered}
\end{equation}
Thus the gradient table is independent of $r_0$, while the ambient
dimension records all three weights.  By
\cref{prop:MR-generator-test}, $\R[R,T,F]$ is a Laplacian algebra.
It is exactly $\cA_R(F)$: writing $S=|\nabla F|^2$, one obtains
\[
 S=4R(aF+bT),\qquad
 \Delta S=8(N+8)(aF+bT)+8(a^2+bh)R^2.
\]
Since $b>0$, the second identity recovers $T$ in $\cA_R(F)$.
The generators are algebraically independent.  Indeed, the nonzero
polynomial $E_0E_1E_2$ is positive somewhere; there the gradients of
the three blocks are independent, and the change to $(R,T,F)$ is
locally invertible.  Fixing $R=1$ therefore leaves spherical quotient
dimension two.

\begin{lemma}[Geometry determined by the quartic data]
\label{lem:quartic-table-geometry}
At each fixed $R\geq0$, the full image of $(T,F)$ is
described by
\begin{equation}\label{eq:quartic-table-image}
 0\leq T\leq R^2/4,\qquad -cT\leq F\leq(1-c)T.
\end{equation}
The critical set of $F$ and the geometric singular set of
$C=\{F=0\}$ are both $Z=\{T=0\}$, and
$\dim Z\leq N-h$.  The hypersurface $C$ is minimal outside $Z$,
both phases approach every point of $C$, and
$\partial\llbracket\{F<0\}\rrbracket$ has support $C$.
\end{lemma}

\begin{proof}
The simplex image in \cref{lem:quartic-block-weights-image} and
homogeneity prove \eqref{eq:quartic-table-image}. Moreover,
\[
 |\nabla F|^2=R\bigl((1-c)^2E_1+c^2E_2\bigr),
\]
so the critical set of $F$ is exactly $Z$.  The reduced table gives
\begin{equation}\label{bh:minimality}
 \cL(F)=-8F\{b(R^2+2T)+2aF\},
\end{equation}
which proves regular minimality.

Nonzero points of $Z$ exist because $E_0$ attains one on the unit
sphere. Fix such a point $p$ and write $R_0=R(p)$ and
$\Hess E_i(p)=8R_0P_i$ for $i=1,2$, using the projections obtained
in the proof of \cref{lem:quartic-block-weights-image}. The quadratic terms of
$\Gamma(E_1,E_2)=0$ give $P_1P_2+P_2P_1=0$.  Taking the trace
shows $P_1P_2=0$.  Hence
\[
 \Hess T(p)=2R_0(P_1+P_2),\qquad
 \Hess F(p)=2R_0\bigl((1-c)P_1-cP_2\bigr).
\]
The first Hessian has rank $h$.  Applying the implicit function
theorem to $h$ independent components of $\nabla T$ places $Z$
locally in a smooth submanifold of codimension $h$.  Consequently
$\dim Z\leq N-h$ and $\cH^{N-1}(Z)=0$.
The second Hessian has positive and negative indices $r_1,r_2$,
respectively, so both phases occur arbitrarily near every such $p$.
If the zero set were a smooth hypersurface at $p$, its leading
quadratic term would vanish on its tangent hyperplane and thus have
rank at most two.  Its actual rank is $h\geq4$, a contradiction.
Thus these points are geometrically singular.

Homogeneity gives the same two-phase conclusion at the origin,
and regularity gives it on $C\setminus Z$.  Polynomial finite
perimeter and the local BV Poincar\'e inequality then identify the
support of the boundary current with $C$.  The vertex is singular
as well: a cone smooth at its vertex would be a hyperplane, whereas
$C$ has nonzero singular points.
\end{proof}

\begin{theorem}[Polynomial subcalibrations from orthogonal quartics]
\label{thm:quartic-table-existence}
For the data \eqref{eq:orthogonal-quartic-data} and
\eqref{eq:quartic-general-weights}--\eqref{eq:quartic-general-generators},
suppose that either
\[
 h\geq9,
 \qquad\text{or}\qquad
 r_1=r_2=4\ \text{ and }\ r_0\geq5.
\]
There exists a homogeneous polynomial $G=AF$ of degree at most ten,
with $A>0$ on $T>0$, whose critical set is $Z=\{T=0\}$ and whose
normalized gradient satisfies
\[
 \operatorname{sgn}(G)\,\diver X_G>0\qquad\text{off }C=\{F=0\}.
\]
The current $\partial\llbracket\{F<0\}\rrbracket$ has full support
$C$ and is strictly area minimizing.  Its regular part is strictly
stable, with the explicit lower bound
\begin{equation}\label{eq:bihomogeneous-Hardy}
 Q_C(\varphi)\geq
 \left[\frac{(N-17)^2}{4}+8h-72\right]
 \int_Cr^{-2}\varphi^2,\qquad
 \varphi\in C_c^\infty(C_{\reg}\setminus\{0\}),\quad r=\sqrt R.
\end{equation}
\end{theorem}

\subsubsection{Selecting the multiplier}

At $p\in Z\setminus\{0\}$, \eqref{bh:minimality} gives
$\lambda_F(p)=-8bR(p)^2<0$. Hence
\cref{lem:critical-multiplier-defect} forces every admissible smooth
multiplier to vanish on $Z$.
For $A\in\R[R,T]$, scaling a nonzero point of $Z$ gives
$A(R,0)=0$ for all $R>0$, so $T$ divides $A$ as a polynomial in
$R,T$.  The lowest two nonzero homogeneous choices are therefore
$T$ and $RT$, up to a positive constant.

These two candidates suffice in the stated ranges:
\begin{equation}\label{bh:potentials}
 G=\begin{cases}
 TF,&h\geq11,\ \text{or }h=10\ \text{and }r_1,r_2\geq3,\\
 RTF,&h=9,\ \text{or }\{r_1,r_2\}=\{2,8\},\\
 RTF,&r_1=r_2=4,\ r_0\geq5.
 \end{cases}
\end{equation}
With $s=T/R^2$ and $z=F/T$, these are $R^4s^2z$ and $R^5s^2z$:
the radial factor changes the defect but preserves vanishing on $Z$.
We use \cref{thm:finite-reduction,thm:exact-multiplier} and, for
$h=9$, the realizability restriction \cref{cor:quartic-nine-weights}.

\subsubsection{The degree-eight potential and its range}

For $G_8=TF$, the reduced table gives
\begin{gather}
 \norm{\nabla G_8}^2=4RT(bT^2+aTF+3F^2),\qquad
 \Delta G_8=R\{2aT+2(h+4)F\},\label{bh:G8-basic}\\
 \begin{split}\label{bh:G8-defect}
 \cL(G_8)=8TF\big\{&bT^2[(h-8)R^2-4T]\\
                  &+aTF[(h-4)R^2-4T]
                    +3F^2[hR^2-4T]\big\}.
 \end{split}
\end{gather}
Introduce the dimensionless variables
\[
 s=T/R^2\in(0,1/4],\qquad z=F/T\in[-c,1-c].
\]
The braces in \eqref{bh:G8-defect}, divided by $R^2T^2$, are
\begin{equation}\label{bh:p8}
 p_8(z,s)=b(h-8-4s)+a(h-4-4s)z+3(h-4s)z^2.
\end{equation}
Since
\[
 b+az=(1-c)^2(z+c)+c^2(1-c-z)>0,
\]
we have $\partial_s p_8=-4(b+az+3z^2)<0$.  It is enough to prove
\begin{equation}\label{bh:q8}
 q(z):=p_8(z,1/4)=b(h-9)+a(h-5)z+3(h-1)z^2>0.
\end{equation}

First take one of the degree-eight cases in \eqref{bh:potentials},
so $h\geq10$. Interchanging $E_1,E_2$ preserves the problem, so we
may assume
$1/(h-2)\leq c\leq1/2$ and $a\geq0$.  For $z\geq0$,
\eqref{bh:q8} is positive.  On $[-c,0]$ its vertex is
$z_*=-a(h-5)/(6(h-1))$, which lies in this interval exactly when
\[
 c\geq c_*:=\frac{h-5}{8(h-2)}.
\]
If $c<c_*$, the minimum is attained at $-c$, where \(q(-c)=4c((h-1)c-1)>0\).
If $c\geq c_*$, the minimum is positive exactly when
\begin{equation}\label{bh:D8}
 D_8(c):=12(h-1)(h-9)c(1-c)-(h-5)^2(1-2c)^2>0.
\end{equation}
This expression increases on $0<c\leq1/2$.  For $h\geq13$,
\[
 D_8(c_*)=
 \frac{3(h-1)(h-5)(h^2-14h+21)}{4(h-2)^2}>0.
\]
For $h=12,11$, the bound $c\geq1/(h-2)>c_*$ gives, respectively,
\[
 D_8(1/10)=107/25>0,\qquad D_8(1/9)=52/27>0.
\]
For $h=10$ with $r_1,r_2\geq3$, one has $c\geq1/4>c_*$ and
$D_8(1/4)=14>0$. This proves positivity in every degree-eight
case of \eqref{bh:potentials}.

The same candidate does not give a uniform result for $h=10$.
For $c=1/8$, hence $a=3/4$ and $b=7/64$, one has
\begin{equation}\label{eq:quartic-degree-eight-obstruction}
 p_8(-1/16,1/4)=-5/256<0.
\end{equation}
By \cref{lem:quartic-table-geometry}, this point is attained in every
realization with $(r_1,r_2)=(2,8)$. Thus $TF$ has the wrong divergence
sign at points of the complement. For $h=9$, one has
$q(z)=4az+24z^2$ with $a\ne0$, so $q$ is negative for small $z$ of
the opposite sign. For $r_1=r_2=4$, one instead has
$q(z)=-b+21z^2$, negative near zero. The exact image assertion again
realizes these negative values with $F\ne0$. Hence, within the
theorem's range, the degree-eight choices in \eqref{bh:potentials}
are precisely the cases in which $TF$ has the required strict sign.

\subsubsection{The degree-ten potential in the remaining ranges}

For $G_{10}=RTF$, the same table gives
\begin{equation}\label{bh:G10}
 \cL(G_{10})=8R^5T^3F\,p_{10}(z,s),
\end{equation}
where
\begin{align}
 p_{10}(z,s)={}&b[h-8+(2h-23)s]
       +a[h-4+(2h-14)s]z\nonumber\\
       &+[3h+(15h-33)s+(18h+27)s^2]z^2\nonumber\\
       &+2(r_0-1)s\bigl[b+az+3(1+3s)z^2\bigr].
       \label{bh:p10}
\end{align}
Because $r_0\geq1$ and $b+az>0$ on the image interval, the last
line is nonnegative.  Thus increasing the unused block weight
$r_0$ can only improve this certificate.

For $h=9$, \cref{cor:quartic-nine-weights} gives $r_0\geq8$.
At $r_0=8$, the full polynomial is
\begin{equation}\label{eq:quartic-nine-certificate}
 p_{10}(z,s)=b(1+9s)+a(5+18s)z+
 (27+144s+315s^2)z^2.
\end{equation}
Since $c\in[1/7,6/7]$, one has $a^2\leq25/49$ and
$4b\geq24/49$. The negative discriminant is therefore bounded below by
\begin{align}
 &4b(1+9s)(27+144s+315s^2)-a^2(5+18s)^2\nonumber\\
 &\qquad\geq
 \frac{23+4788s+30564s^2+68040s^3}{49}>0.
 \label{eq:quartic-nine-discriminant}
\end{align}
The leading coefficient is positive, so \eqref{eq:quartic-nine-certificate}
is positive for all real $z$ and $0<s\leq1/4$. Increasing $r_0$
preserves positivity on the image interval.

When $h=10$, the first three terms are
\[
 b(2-3s)+a(6+6s)z+(30+117s+207s^2)z^2.
\]
Here $c\in[1/8,7/8]$, so $a^2\leq9/16$ and
$b=(1-a^2)/4$.  The negative of the discriminant is bounded below by
\begin{align*}
 &4b(2-3s)(30+117s+207s^2)-a^2(6+6s)^2\\
 &\hspace{10mm}\geq
 6+\frac{45}{2}s+\frac{117}{16}s^2-\frac{4347}{16}s^3
 \geq\frac{1797}{1024}>0.
\end{align*}
The final inequality follows from $0<s\leq1/4$ by discarding the
positive linear and quadratic terms. The leading coefficient is
positive, so this polynomial is positive for all real $z$ and all
allowed $s$.

Finally, for $r_1=r_2=4$ one has $h=8$, $a=0$, and $b=1/4$.
The full polynomial becomes
\begin{equation}\label{eq:quartic-balanced-eight-certificate}
 p_{10}(z,s)=\frac{2r_0-9}{4}s+
 \bigl[24+(81+6r_0)s+(153+18r_0)s^2\bigr]z^2.
\end{equation}
This is positive for every $s>0$ if $r_0\geq5$.
Together these calculations prove the required sign in every
case of \eqref{bh:potentials}.

\subsubsection{Global minimality and the Hardy bound}

\begin{proof}[Completion of the proof of \cref{thm:quartic-table-existence}]
\leavevmode\par\nobreak
Both multipliers are positive on $T>0$. Euler's identity excludes
critical points where $F\ne0$, and $|\nabla F|^2=4bRT$ excludes them
on $F=0$, $T>0$. On $Z$, both potentials and their gradients vanish.
Thus they have critical set $Z$ and the same phases as $F$.
The full-support and dimension conclusions of
\cref{lem:quartic-table-geometry} verify the hypotheses of
\cref{cor:gradient-subcalibration,thm:removable-subcalibration}.
Their comparison, followed by \cref{cor:polynomial-strict-minimality},
gives strict integral-current mass minimization and the gap
$\Theta\varepsilon^{N-1}$.

For stability, \eqref{eq:boundary-potential-algebraic} gives
\[
 \beta_{TF}=2(h-8)\frac{R^2}{T}-8,\qquad
 \beta_{RTF}=2(h-8)\frac{R^2}{T}+4\rho-50.
\]
Since $R^2/T\geq4$ and $h\geq8$, the homogeneous Hardy estimate
\cref{thm:ground-state-gap}, with cone dimension $N-1=2\rho-3$,
gives the same lower bound for both potential degrees:
\begin{align*}
 \left(8-\frac{N-1}{2}\right)^2+8(h-8)-8
 &=\frac{(N-17)^2}{4}+8h-72,\\
 \left(10-\frac{N-1}{2}\right)^2+8(h-8)+4\rho-50
 &=\frac{(N-17)^2}{4}+8h-72.
\end{align*}
For $h\geq10$ this coefficient is positive. For $h=9$,
\cref{cor:quartic-nine-weights} gives $N=2r_0+16\geq32$, so the
coefficient $(N-17)^2/4$ is positive as well. In the remaining
case, $h=8$ and $r_0\geq5$ imply $N\geq24$, so it is at least
$17/4$.  This proves \eqref{eq:bihomogeneous-Hardy}.
\end{proof}

For the splitting specialization $r_0=1$, $h=\ell$, one has
$N=2\ell$, and the Hardy coefficient becomes
$\ell^2-9\ell+1/4$. For $r_0>1$ and unequal active weights,
\cref{prop:quartic-weight-realizability} determines realizability,
and the inequalities above give a sufficient minimizing range.

\subsubsection{The balanced specialization}

\begin{corollary}[The degree-eight potential for balanced quartics]
\label{prop:balanced-quartic-multiplier}
Suppose that the three nonzero quartics in
\eqref{eq:orthogonal-quartic-data} have weights $(r_0,\nu,\nu)$,
with $\nu\geq5$. Set
\[
 F=\frac{E_+-E_-}{4},\qquad W=\frac{E_++E_-}{4},\qquad G=WF.
\]
The conclusions of \cref{thm:quartic-table-existence} hold for $G$,
with $\Sigma=\{W=0\}$. In this normalization, $\Delta F=0$,
$|F|\leq W\leq R^2/4$, and
\begin{equation}\label{eq:balanced-quartic-smallest-algebra}
 \cA_R(F)=\R[R,W,F]
\end{equation}
is a polynomial algebra with generator degrees $(2,4,4)$.
Moreover,
\begin{align}
 |\nabla G|^2&=4RW(W^2+3F^2),\qquad
 \Delta G=4(\nu+2)RF,\label{eq:balanced-quartic-gradient}\\
 \cL(G)&=16FW\left\{
 W^2\big[(\nu-4)R^2-2W\big]
 +3F^2\big[\nu R^2-2W\big]\right\}.
 \label{eq:balanced-quartic-certificate}
\end{align}
\end{corollary}

\begin{proof}
Take $T=W$, $F_{\mathrm{gen}}=F/2$, $h=2\nu\geq10$, and $c=1/2$
in \cref{thm:quartic-table-existence}. Its table becomes
\begin{equation}\label{eq:balanced-quartic-table}
 \Delta(R,W,F)=(2N,4\nu R,0),\qquad
 (\Gamma(I_i,I_j))_{I=(R,W,F)}=
 \begin{pmatrix}
 4R&8W&8F\\
 8W&4RW&4RF\\
 8F&4RF&4RW
 \end{pmatrix}.
\end{equation}
Substituting this normalization in
\eqref{bh:G8-basic}--\eqref{bh:G8-defect} gives
\eqref{eq:balanced-quartic-gradient}--\eqref{eq:balanced-quartic-certificate}.
\end{proof}

The following table organizes the applications of the common
quartic theorem. The splitting, Riedler, and weight-two families
have exact thresholds, obtained by combining this theorem with
instability arguments. Each construction also identifies its
singular set, including the real rank-one strata in the weight-two
families.

\begin{table}[htbp]
\centering\small
\setlength{\tabcolsep}{5pt}
\renewcommand{\arraystretch}{1.15}
\caption{Euclidean quartic realizations of
\cref{thm:quartic-table-existence}. Only parameters admitting the
specified Clifford realization are included. The intrinsic
Simons-cone result in \cref{thm:Hopf-Simons-section} uses a
separate differential table.}
\label{tab:quartic-weight-realizations}
\begin{tabular}{@{}lll@{}}
\toprule
Family & $(r_0,r_1,r_2)$ & Minimizing parameters used here\\
\midrule
OT--FKM splitting & $(1,m,\ell-m)$ & $\ell\ge10$, $2\le m\le\ell-2$\\
Riedler pullback & $(N/2-2k+1,k,k)$ & $N\ge32$, $k\ge4$\\
Hopf pairing & $(1,8,8)$ & Ambient dimension $32$\\
Balanced Clifford example & $(3,7,7)$ & Ambient dimension $32$\\
Additional block example & $(2,5,10)$ & Ambient dimension $32$\\
Weight-two families & $(2,m,4k-m-1)$ & $m=2,3$, $k\ge3$\\
Unequal Lawson pullback & $(22,5,6)$ & Ambient dimension $64$\\
\bottomrule
\end{tabular}
\end{table}
